\documentclass[../../main.tex]{subfiles}
\begin{document}

\begin{center}
  {\fontsize{48pt}{58pt}\selectfont \textbf{2000}\normalsize\textbf{年度}}
\end{center}

\begin{table}[htb]
  \centering
  \caption{出題テーマと難易度}
  \begin{tabular}{|c|p{9cm}|c|}
    \hline
    \textbf{問題番号} & \textbf{テーマ} & \textbf{難易度} \\
    \hline
    第1問 & 円に内接する四角形のベクトル & \\
    \hline
    第2問 & 平方根のグラフと複素数解の絶対値 & \\
    \hline
    第3問 & 空間ベクトルの内積と角度領域 & \\
    \hline
    第4問 & 複素整数の素数乗 & \\
    \hline
    第5問 & 余弦の定積分で定まる数列 & \\
    \hline
    第6問 & サイコロの和の剰余と極限 & \\
    \hline
  \end{tabular}
\end{table}

\begin{table}[htb]
  \centering
  \caption{解答}
  \renewcommand{\arraystretch}{1.4}
  \begin{tabular}{|c|c|p{9cm}|}
    \hline
    \textbf{問題番号} & \textbf{小問} & \textbf{答え} \\
    \hline
    第1問 & -- & $\overrightarrow{AP}=\dfrac{(1-p)\overrightarrow{AB}+p\overrightarrow{AC}}{p^2-p+1}$ \\
    \hline
    第2問 & (1) & $0<a\le4-2\sqrt{2}$ \\
    \cline{2-3}
    & (2) & $\min|\beta|=\dfrac{3-2\sqrt{2}}{2}$ \\
    \hline
    第3問 & (1) & 指定された内積の不等式が成立（証明問題） \\
    \cline{2-3}
    & (2) & $0\le\alpha,\beta\le\pi$で$\cos^2\alpha-\cos\alpha\cos\beta+\cos^2\beta\le3/4$の領域 \\
    \hline
    第4問 & -- & $(a+bi)^p$は実数ではない（証明問題） \\
    \hline
    第5問 & (1) & $c_{n+2}=-\dfrac{(n+3)(n+2)}{\pi^2}(1+c_n)$ \\
    \cline{2-3}
    & (2) & $\lim_{n\to\infty}c_n=-1$ \\
    \cline{2-3}
    & (3) & $\lim_{n\to\infty}\dfrac{c_{n+1}+1}{c_n+1}=1$ \\
    \hline
    第6問 & (1) & $p_{n+1}(k)=\dfrac{1+p_n(k-1)}{6}$（添字は$5$を法として読む） \\
    \cline{2-3}
    & (2) & $m_n\le1/5\le M_n$，$M_{n+1}-m_{n+1}\le(M_n-m_n)/6$ \\
    \cline{2-3}
    & (3) & $\lim_{n\to\infty}p_n(k)=1/5$ \\
    \hline
  \end{tabular}
\end{table}

\paragraph{試験形態}
6問．試験時間は要確認．

\end{document}
